\documentclass[12pt,a4paper]{article}
\usepackage[utf8]{inputenc}
\usepackage[T1]{fontenc}
\usepackage[english]{babel}
\usepackage{geometry}
\geometry{margin=2.5cm}
\usepackage{graphicx}
\usepackage{booktabs}
\usepackage{hyperref}
\usepackage{listings}
\usepackage{xcolor}
\usepackage{amsmath}
\usepackage{float}
\usepackage{enumitem}
\usepackage{fancyhdr}
\usepackage{titlesec}

\definecolor{codeblue}{RGB}{0,102,204}
\definecolor{codegray}{RGB}{128,128,128}
\definecolor{codegreen}{RGB}{0,128,0}
\definecolor{backcolour}{RGB}{245,245,245}

\lstdefinestyle{mystyle}{
    backgroundcolor=\color{backcolour},
    commentstyle=\color{codegreen},
    keywordstyle=\color{codeblue},
    numberstyle=\tiny\color{codegray},
    basicstyle=\ttfamily\small,
    breaklines=true,
    frame=single,
    numbers=left,
    numbersep=5pt,
}
\lstset{style=mystyle}

\pagestyle{fancy}
\fancyhf{}
\rhead{Smart Document Analyst}
\lhead{UIR -- S8 Integrated Project}
\rfoot{Page \thepage}

\title{
    \vspace{-1cm}
    \textbf{Smart Document Analyst}\\[0.3cm]
    \large A Multi-Agent AI System for Document Classification,\\
    Extraction, Summarization, and Report Generation\\[0.5cm]
    \normalsize S8 Integrated Project -- Building Multi-Agent AI Systems\\
    AI \& Big Data Program -- UIR 2025--2026
}
\author{
    Benmouma Salma \and Gassi Oumaima\\[0.3cm]
    \small Supervised by Prof.\ Hakim Hafidi
}
\date{April 2026}

\begin{document}
\maketitle
\thispagestyle{empty}

\begin{abstract}
This report presents the design, implementation, and evaluation of a multi-agent AI system for automated document analysis. The system employs three specialist agents orchestrated by CrewAI to classify documents using a CNN deep learning model, extract and summarize text content, and generate structured analysis reports. A fine-tuned ResNet-18 model trained on the RVL-CDIP dataset achieves 85.62\% test accuracy across 8 document classes. The system includes a human-in-the-loop checkpoint, comprehensive error handling, and JSON-based action logging. This report details the architectural decisions, model training, agent design, and evaluation results.
\end{abstract}

\tableofcontents
\newpage

% ============================================================
\section{Introduction}
% ============================================================

\subsection{Context and Motivation}
In modern organizations, the volume of documents---invoices, letters, reports, forms---continues to grow rapidly. Manually classifying, reading, and summarizing these documents is time-consuming and error-prone. An intelligent system capable of automating this pipeline would significantly improve productivity and accuracy.

\subsection{Project Objectives}
The objective of this project is to design and build a \textbf{multi-agent AI system} where specialized agents collaborate to:
\begin{enumerate}
    \item \textbf{Classify} documents into predefined categories using a CNN deep learning model.
    \item \textbf{Extract} text from documents (native PDF or OCR for scanned documents).
    \item \textbf{Summarize} the extracted content using a large language model (Gemini API).
    \item \textbf{Generate} a structured professional report combining all analysis results.
\end{enumerate}

\subsection{Technical Constraints}
As specified in the project brief, the system must satisfy the following requirements:
\begin{itemize}
    \item At least 2 specialist agents + 1 orchestrator
    \item A PyTorch deep learning model trained by the team, integrated as a functional tool
    \item At least 2 tools with clear input/output schemas
    \item A human-in-the-loop checkpoint requiring human approval
    \item Error handling for tool failures, invalid inputs, and unexpected outputs
    \item JSON logging with timestamps for every agent action
\end{itemize}

% ============================================================
\section{System Architecture}
% ============================================================

\subsection{High-Level Overview}
The system follows a \textbf{sequential pipeline} architecture where documents flow through four stages: classification, human approval, extraction \& summarization, and report generation. Each stage is handled by a specialist agent equipped with dedicated tools.

\begin{figure}[H]
\centering
\begin{verbatim}
    User Input (PDF/Image)
         |
    [Orchestrator Agent] --- Manages workflow
         |
    [Classifier Agent] --- CNN Model Tool (PyTorch)
         |
    [HITL Checkpoint] --- Human approves/corrects classification
         |
    [Extractor & Summarizer Agent] --- OCR + LLM Tools
         |
    [Report Generator Agent] --- Report Builder Tool
         |
    Final Analysis Report (Markdown/PDF)
\end{verbatim}
\caption{System pipeline architecture}
\end{figure}

\subsection{Agent Roles and Responsibilities}

\begin{table}[H]
\centering
\caption{Agent definitions and their tools}
\begin{tabular}{@{}llp{5.5cm}@{}}
\toprule
\textbf{Agent} & \textbf{Tools} & \textbf{Responsibility} \\
\midrule
Classifier & \texttt{cnn\_classify\_tool} & Classify document type using the trained CNN model \\
Extractor \& Summarizer & \texttt{ocr\_extract\_tool}, \texttt{llm\_summarize\_tool} & Extract text via OCR and generate a type-aware summary \\
Report Generator & \texttt{report\_builder\_tool} & Produce a structured Markdown/PDF report \\
\bottomrule
\end{tabular}
\end{table}

\subsection{Orchestration with CrewAI}
The agents are orchestrated using \textbf{CrewAI} with a sequential process. The orchestrator:
\begin{itemize}
    \item Routes tasks to agents in the correct order
    \item Passes context (classification results) between agents
    \item Enforces the HITL checkpoint between classification and extraction
    \item Aggregates the final output
\end{itemize}

The system supports two execution modes:
\begin{itemize}
    \item \textbf{Standard mode}: Fully automated sequential pipeline via \texttt{crew.run()}
    \item \textbf{HITL mode}: Manual orchestration with an interactive classification approval checkpoint via \texttt{crew.run\_with\_hitl()}
\end{itemize}

\subsection{Project Structure}
\begin{lstlisting}[language=bash, caption=Repository structure]
smart-document-analyst/
├── config/
│   ├── agents.yaml          # Agent role definitions
│   └── tasks.yaml           # Task pipeline definitions
├── src/
│   ├── main.py              # CLI entry point
│   ├── crew.py              # CrewAI crew orchestration
│   ├── agents/              # Agent definitions
│   ├── tools/               # Tool implementations
│   ├── models/              # PyTorch model definition
│   └── utils/               # Logger, HITL, preprocessing
├── model/
│   └── document_classifier.pt  # Trained model weights
├── notebooks/
│   └── train_on_colab.ipynb    # Training notebook
└── tests/                      # Unit tests
\end{lstlisting}

% ============================================================
\section{Deep Learning Model}
% ============================================================

\subsection{Dataset: RVL-CDIP}
We use the \textbf{RVL-CDIP} (Ryerson Vision Lab -- Complex Document Information Processing) dataset, a widely-used benchmark for document image classification. From the full 16-class dataset, we selected a \textbf{subset of 8 classes} to maintain feasibility on Google Colab's free tier:

\begin{table}[H]
\centering
\caption{Selected document classes (2,000 images per class)}
\begin{tabular}{@{}ll@{}}
\toprule
\textbf{Class} & \textbf{Description} \\
\midrule
letter & Formal correspondence \\
form & Structured forms with fields \\
invoice & Financial billing documents \\
resume & Curriculum vitae / CVs \\
scientific\_publication & Academic papers \\
email & Electronic mail printouts \\
memo & Internal memoranda \\
advertisement & Marketing materials \\
\bottomrule
\end{tabular}
\end{table}

Total dataset size: \textbf{16,000 images}, split as follows:
\begin{itemize}
    \item Training: 12,800 images (80\%)
    \item Validation: 1,600 images (10\%)
    \item Test: 1,600 images (10\%)
\end{itemize}

\subsection{Model Architecture}
We use a \textbf{fine-tuned ResNet-18} pretrained on ImageNet. The architecture choices are:

\begin{itemize}
    \item \textbf{Backbone}: ResNet-18 with ImageNet pretrained weights
    \item \textbf{Frozen layers}: \texttt{conv1}, \texttt{bn1}, \texttt{layer1}, \texttt{layer2} (early feature extraction layers)
    \item \textbf{Trainable layers}: \texttt{layer3}, \texttt{layer4}, and the custom classification head
    \item \textbf{Classification head}: Dropout(0.3) $\rightarrow$ Linear(512, 256) $\rightarrow$ ReLU $\rightarrow$ BatchNorm $\rightarrow$ Dropout(0.15) $\rightarrow$ Linear(256, 8)
\end{itemize}

\textbf{Rationale}: Fine-tuning a pretrained model leverages learned visual features while adapting to document-specific patterns. Freezing early layers reduces training time and prevents overfitting on our relatively small dataset.

\textbf{Parameters}: 11,310,408 total; trainable parameters focus on the deeper layers and classification head.

\subsection{Training Configuration}

\begin{table}[H]
\centering
\caption{Training hyperparameters}
\begin{tabular}{@{}ll@{}}
\toprule
\textbf{Parameter} & \textbf{Value} \\
\midrule
Input size & 224 $\times$ 224 (grayscale $\rightarrow$ 3-channel) \\
Batch size & 32 \\
Optimizer & AdamW (lr=$10^{-4}$, weight\_decay=$10^{-4}$) \\
Scheduler & CosineAnnealingLR ($T_{max}=20$, $\eta_{min}=10^{-6}$) \\
Loss function & CrossEntropyLoss \\
Epochs & 20 (early stopping, patience=5) \\
Device & NVIDIA T4 GPU (Google Colab) \\
\bottomrule
\end{tabular}
\end{table}

\textbf{Data augmentation} (training only): RandomCrop(224), RandomRotation(5\textdegree), RandomAffine (translate=5\%, scale=95--105\%), ColorJitter (brightness=0.2, contrast=0.2).

\subsection{Training Results}
Training converged in \textbf{18 epochs} before early stopping was triggered. The best validation accuracy was achieved at epoch 13.

\begin{table}[H]
\centering
\caption{Training progression (selected epochs)}
\begin{tabular}{@{}ccccc@{}}
\toprule
\textbf{Epoch} & \textbf{Train Loss} & \textbf{Train Acc} & \textbf{Val Loss} & \textbf{Val Acc} \\
\midrule
1  & 0.9478 & 67.51\% & 0.7341 & 74.62\% \\
5  & 0.3796 & 87.30\% & 0.5043 & 83.81\% \\
10 & 0.1881 & 93.80\% & 0.5024 & 86.19\% \\
13 & 0.1231 & 95.78\% & 0.5075 & \textbf{87.62\%} \\
18 & 0.0693 & 98.02\% & 0.5397 & 87.62\% \\
\bottomrule
\end{tabular}
\end{table}

\subsection{Test Set Evaluation}

\begin{table}[H]
\centering
\caption{Per-class classification results on the test set}
\begin{tabular}{@{}lcccc@{}}
\toprule
\textbf{Class} & \textbf{Precision} & \textbf{Recall} & \textbf{F1-Score} & \textbf{Support} \\
\midrule
letter & 0.88 & 0.95 & 0.92 & 200 \\
form & 0.91 & 0.95 & 0.93 & 191 \\
invoice & 0.78 & 0.72 & 0.75 & 207 \\
resume & 0.82 & 0.83 & 0.83 & 204 \\
scientific\_publication & 0.83 & 0.71 & 0.76 & 201 \\
email & 0.81 & 0.84 & 0.82 & 206 \\
memo & 0.89 & 0.90 & 0.90 & 196 \\
advertisement & 0.92 & 0.94 & 0.93 & 195 \\
\midrule
\textbf{Overall Accuracy} & \multicolumn{4}{c}{\textbf{85.62\%}} \\
\textbf{Macro Avg} & 0.86 & 0.86 & 0.86 & 1600 \\
\bottomrule
\end{tabular}
\end{table}

\textbf{Analysis}: The model performs best on advertisement (F1=0.93) and form (F1=0.93), and weakest on invoice (F1=0.75) and scientific\_publication (F1=0.76). The confusion between invoice and form is expected given their visual similarity (both contain structured fields and tables).

% ============================================================
\section{Tools Implementation}
% ============================================================

Each tool is implemented as a CrewAI \texttt{BaseTool} subclass with a Pydantic input schema.

\subsection{Tool 1: CNN Classification Tool}
\begin{itemize}
    \item \textbf{Input}: File path to a document (PDF, PNG, JPG)
    \item \textbf{Output}: JSON with predicted class, confidence score, and top-3 predictions
    \item \textbf{Process}: Loads the trained PyTorch model, preprocesses the document image (PDF pages are rendered via PyMuPDF), runs inference with \texttt{torch.no\_grad()}, and applies softmax for probabilities
\end{itemize}

\subsection{Tool 2: OCR Extraction Tool}
\begin{itemize}
    \item \textbf{Input}: File path to a document
    \item \textbf{Output}: JSON with extracted text, page count, and extraction method
    \item \textbf{Strategy}: Tries PyMuPDF native text extraction first; if the result is too short ($<$50 characters, indicating a scanned document), falls back to Tesseract OCR
\end{itemize}

\subsection{Tool 3: LLM Summarization Tool}
\begin{itemize}
    \item \textbf{Input}: Extracted text, document type, maximum summary length
    \item \textbf{Output}: JSON with summary, key points, and word count
    \item \textbf{Features}: Adapts the prompt based on document type (e.g., invoices focus on amounts; resumes focus on skills). Includes retry logic with exponential backoff and an extractive fallback when the API is unavailable
\end{itemize}

\subsection{Tool 4: Report Builder Tool}
\begin{itemize}
    \item \textbf{Input}: Classification, extraction, and summary results as JSON
    \item \textbf{Output}: JSON with the path to the generated report file
    \item \textbf{Format}: Generates a professional Markdown report with optional PDF conversion via \texttt{fpdf2}
\end{itemize}

% ============================================================
\section{Human-in-the-Loop}
% ============================================================

The HITL checkpoint is triggered after classification and before downstream processing. It presents the CNN's prediction to the user and offers three options:

\begin{enumerate}
    \item \textbf{Approve}: Accept the classification and proceed
    \item \textbf{Override}: Correct the class label (the corrected label propagates to the summarizer and report)
    \item \textbf{Reject}: Stop the pipeline entirely for this document
\end{enumerate}

This is not decorative---the user's decision has genuine functional impact on the system's behavior. If the user overrides the class, the Extractor Agent adapts its summarization style to the corrected document type, producing a more relevant summary.

% ============================================================
\section{Error Handling and Logging}
% ============================================================

\subsection{Error Handling}
Every tool implements defensive error handling:

\begin{table}[H]
\centering
\caption{Error handling scenarios}
\begin{tabular}{@{}lp{7cm}@{}}
\toprule
\textbf{Scenario} & \textbf{Handling} \\
\midrule
File not found & Returns descriptive error JSON, no crash \\
Unsupported format & Validates extension, returns supported formats list \\
CNN model missing & Graceful error with download instructions \\
OCR returns empty & Falls back to Tesseract; if still empty, warns user \\
Gemini API failure & 3 retries with exponential backoff, then extractive fallback \\
Invalid LLM output & Parse validation and re-prompting \\
\bottomrule
\end{tabular}
\end{table}

\subsection{Logging}
All agent actions are logged to \texttt{logs/agent\_actions.jsonl} in structured JSON format:

\begin{lstlisting}[language=Python, caption=Example log entry]
{
  "timestamp": "2026-04-26T14:23:01Z",
  "agent": "ClassifierAgent",
  "action": "cnn_classify_tool",
  "input": {"file_path": "data/sample_docs/invoice.pdf"},
  "output": {"class": "invoice", "confidence": 0.943},
  "status": "success",
  "duration_ms": 1230
}
\end{lstlisting}

% ============================================================
\section{Evaluation and Robustness}
% ============================================================

\subsection{CNN Model Evaluation}
The CNN model was rigorously evaluated with:
\begin{itemize}
    \item \textbf{Accuracy}: 85.62\% on a held-out test set of 1,600 images
    \item \textbf{Confusion matrix}: Analyzed to identify class-pair confusions
    \item \textbf{Per-class metrics}: Precision, recall, and F1-score for all 8 classes
    \item \textbf{Training curves}: Monitored for overfitting (early stopping triggered at epoch 18)
\end{itemize}

\subsection{Honest Failure Analysis}
\begin{itemize}
    \item \textbf{Invoice vs.\ Form confusion}: Both document types share visual features (tables, structured fields), leading to the lowest F1 scores
    \item \textbf{Scientific publication recall}: At 71\%, some academic papers are misclassified as letters or memos, likely due to text-heavy layouts that resemble correspondence
    \item \textbf{Overfitting gap}: Training accuracy (98\%) vs.\ test accuracy (85.6\%) indicates some overfitting despite data augmentation and dropout. More training data or stronger regularization could help
\end{itemize}

% ============================================================
\section{Challenges and Lessons Learned}
% ============================================================

\begin{enumerate}
    \item \textbf{Dependency conflicts on Colab}: NumPy 2.0 broke compatibility with scikit-learn and scipy. Solution: pinning \texttt{numpy<2.0.0} and restarting the runtime.
    \item \textbf{HuggingFace dataset loading}: The \texttt{rvl\_cdip} shorthand was deprecated. Solution: using the full path \texttt{aharley/rvl\_cdip}.
    \item \textbf{GitHub notebook rendering}: Jupyter widget metadata from tqdm caused ``Invalid Notebook'' errors. Solution: stripping \texttt{metadata.widgets} before committing.
    \item \textbf{Agent design}: The key insight is that each agent must have a \textit{distinct, justified role}. Merging tools into fewer agents would reduce the system's modularity and make it harder to debug.
\end{enumerate}

% ============================================================
\section{Conclusion}
% ============================================================

We have designed, built, and evaluated a multi-agent AI system for document analysis that satisfies all project requirements. The system demonstrates genuine collaboration between specialized agents, each with a distinct role and dedicated tools. The CNN deep learning model, trained on real document data, serves as a functional tool---not decoration---within the agent pipeline.

The system achieves 85.62\% classification accuracy, includes a meaningful human-in-the-loop checkpoint, and handles errors gracefully without crashing. Every agent action is logged with timestamps in structured JSON format.

\textbf{Future work} could include:
\begin{itemize}
    \item Training on the full 16-class RVL-CDIP dataset for broader coverage
    \item Adding a Streamlit/Gradio web interface for non-technical users
    \item Migrating to LangGraph for more complex agent communication patterns
    \item Fine-tuning a smaller language model locally to remove API dependency
\end{itemize}

% ============================================================
\section*{References}
% ============================================================
\begin{enumerate}
    \item Harley, A.W., Ufkes, A., Derpanis, K.G. (2015). \textit{Evaluation of Deep Convolutional Nets for Document Image Classification and Retrieval}. ICDAR.
    \item He, K. et al.\ (2016). \textit{Deep Residual Learning for Image Recognition}. CVPR.
    \item CrewAI Documentation. \url{https://docs.crewai.com/}
    \item PyTorch Documentation. \url{https://pytorch.org/docs/}
    \item Google Gemini API. \url{https://ai.google.dev/}
    \item RVL-CDIP Dataset. \url{https://huggingface.co/datasets/aharley/rvl_cdip}
\end{enumerate}

\end{document}
